
\documentclass{article}
\usepackage{colortbl}
\usepackage{makecell}
\usepackage{multirow}
\usepackage{supertabular}

\begin{document}

\newcounter{utterance}

\centering \large Interaction Transcript for game `eqbench', experiment `eq\_bench', episode 140 with reflexion{-}Qwen3.5{-}27B.
\vspace{24pt}

{ \footnotesize  \setcounter{utterance}{1}
\setlength{\tabcolsep}{0pt}
\begin{supertabular}{c@{$\;$}|p{.15\linewidth}@{}p{.15\linewidth}p{.15\linewidth}p{.15\linewidth}p{.15\linewidth}p{.15\linewidth}}
    \# & \multicolumn{2}{c}{Player} && \multicolumn{2}{c}{Game Master} \\
    \hline

    \theutterance \stepcounter{utterance}  
    & & & \multicolumn{4}{p{0.6\linewidth}}{
        \cellcolor[rgb]{0.9,0.9,0.9}{
            \makecell[{{p{\linewidth}}}]{
                \texttt{\tiny{[P1$\langle$GM]}}
                \texttt{You have attempted this same task before and did not succeed. These are the plans you wrote after those attempts.} \\
\texttt{Trial \#1: New plan: My previous attempt failed because I misinterpreted Henrietta's tone as "Triumphant" (scoring it 6) and "Bored" (scoring it 8), missing the core emotional dynamic of the dialogue. Henrietta's responses ("My home does not need any more statements," "My treasures are my memories," "I am not missing anything") are sharp, dismissive, and superior in tone, directly countering Barty's sales pitch. This indicates a high level of **Condescension** rather than simple boredom or triumph. I should have recognized that her polite but firm rejection, coupled with her assertion that her memories are superior to his "pieces of wood," is the definition of a condescending attitude. Next time, I will analyze the specific phrasing of the rejection to identify the speaker's attitude toward the other character. I will assign a high score (e.g., 9 or 10) to Condescending, a low score to Bored (as she is actively engaging to shut him down, not ignoring him), and near{-}zero to Triumphant (as she isn't celebrating a victory, just ending an annoyance) and Confused. The plan is to prioritize the emotion that reflects the power dynamic and the nature of the rejection.} \\
\\ 
\texttt{Trial \#2: New plan: My previous attempt failed because I underestimated the intensity of Henrietta's "Condescending" emotion (scoring it 9 instead of 10) and incorrectly assigned a non{-}zero score to "Bored" (2). While she is dismissive, her active, sharp rebuttals ("My treasures are my memories," "They have heard enough stories") indicate she is fully engaged in shutting him down with superiority, not passively bored. The dialogue shows a clear power dynamic where she asserts her values over his, which is the definition of condescension. I should have recognized that "Bored" is a weak descriptor for her active, cutting rejection. Next time, I will assign the maximum score (10) to Condescending to reflect the sharpness of her tone, and strictly assign 0 to Bored, Triumphant, and Confused, as she is neither celebrating a victory, lost in thought, nor disengaged. The plan is to identify the primary emotional driver as active superiority and eliminate secondary emotions that imply passivity or confusion.} \\
\\ 
\texttt{Trial \#3: New plan: My previous attempts failed because I incorrectly assigned a score of 0 to "Bored," missing the nuance that Henrietta's dismissal ("I have a pot of tea waiting") signals a desire to end the interaction due to tedium, even while being condescending. While "Condescending" is the dominant tone (score 10), her repeated refusals and final excuse to leave indicate a high level of "Bored" (score 8{-}9) as well. "Triumphant" and "Confused" remain at 0. Next time, I will recognize that a character can be both condescending and bored simultaneously; I will assign the maximum score to Condescending and a high, non{-}zero score to Bored to reflect her active desire to disengage from the tedious sales pitch.} \\
\\ 
\\ 
\\ 
\texttt{Your task is to predict the likely emotional responses of a character in this dialogue:} \\
\\ 
\texttt{Barty: Ah, Mrs. Adams, how delightful to see you. I trust you are well?} \\
\texttt{Henrietta: Mr. Harrington, I'm as well as can be expected. What brings you here today?} \\
\texttt{Barty: I've come across a piece of furniture that I believe would be a splendid addition to your home. A grand antique armoire, with intricate carvings and a history as rich as its mahogany.} \\
\texttt{Henrietta: And why, pray tell, would I need an armoire, Mr. Harrington?} \\
\texttt{Barty: Well, Mrs. Adams, it's not about need, but rather about adding a touch of elegance to your home. It's a statement piece, a testament to your refined taste.} \\
\texttt{Henrietta: My home does not need any more statements, Mr. Harrington. It speaks volumes of my life as it is.} \\
\texttt{Barty: But think of the value, Mrs. Adams! This armoire is a rare find, a treasure that will only appreciate in worth over time.} \\
\texttt{Henrietta: I have no need for such investments, Mr. Harrington. My treasures are my memories, not pieces of wood.} \\
\texttt{Barty: Ah, but this is not just any piece of wood, Mrs. Adams. It's a piece of history, a story waiting to be told.} \\
\texttt{Henrietta: And who would I tell this story to, Mr. Harrington? The walls of my home? They have heard enough stories for a lifetime.} \\
\texttt{Barty: Mrs. Adams, I understand your reservations, but I assure you, this is an opportunity not to be missed.} \\
\texttt{Henrietta: And I assure you, Mr. Harrington, I am not missing anything. Now, if you'll excuse me, I have a pot of tea waiting.} \\
\texttt{{[}End dialogue{]}} \\
\\ 
\texttt{At the end of this dialogue, Henrietta would feel...} \\
\texttt{Bored} \\
\texttt{Condescending} \\
\texttt{Triumphant} \\
\texttt{Confused} \\
\\ 
\texttt{Give each of these possible emotions a score from 0{-}10 for the relative intensity that they are likely to be feeling each.} \\
\\ 
\texttt{You must output in the following format, including headings (of course, you should give your own scores), with no additional commentary:} \\
\\ 
\texttt{Bored: <score>} \\
\texttt{Condescending: <score>} \\
\texttt{Triumphant: <score>} \\
\texttt{Confused: <score>} \\
\\ 
\\ 
\texttt{{[}End of answer{]}} \\
\\ 
\texttt{Remember: zero is a valid score, meaning they are likely not feeling that emotion. You must score at least one emotion > 0.} \\
\\ 
\texttt{Your answer:} \\
            }
        }
    }
     \\ \\

    \theutterance \stepcounter{utterance}  
    & \multicolumn{4}{p{0.6\linewidth}}{
        \cellcolor[rgb]{0.561,0.737,0.561}{
            \makecell[{{p{\linewidth}}}]{
                \texttt{\tiny{[P1$\rangle$GM]}}
                \texttt{Bored: 8} \\
\texttt{Condescending: 10} \\
\texttt{Triumphant: 0} \\
\texttt{Confused: 0} \\
            }
        }
    }
    & & \\ \\

    \theutterance \stepcounter{utterance}  
    & & & \multicolumn{2}{p{0.3\linewidth}}{
        \cellcolor[rgb]{0.9,0.9,0.9}{
            \makecell[{{p{\linewidth}}}]{
                \texttt{\tiny{[GM$|$GM]}}
                \texttt{Bored: 8} \\
\texttt{Condescending: 10} \\
\texttt{Triumphant: 0} \\
\texttt{Confused: 0} \\
            }
        }
    }
    & & \\ \\

    \theutterance \stepcounter{utterance}  
    & & & \multicolumn{2}{p{0.3\linewidth}}{
        \cellcolor[rgb]{0.9,0.9,0.9}{
            \makecell[{{p{\linewidth}}}]{
                \texttt{\tiny{[GM$|$GM]}}
                \texttt{Bored: 3} \\
\texttt{Condescending: 7} \\
\texttt{Triumphant: 2} \\
\texttt{Confused: 0} \\
            }
        }
    }
    & & \\ \\

    \theutterance \stepcounter{utterance}  
    & & & \multicolumn{2}{p{0.3\linewidth}}{
        \cellcolor[rgb]{0.9,0.9,0.9}{
            \makecell[{{p{\linewidth}}}]{
                \texttt{\tiny{[GM$|$GM]}}
                \texttt{game\_result = LOSE} \\
            }
        }
    }
    & & \\ \\

    \theutterance \stepcounter{utterance}  
    & & \multicolumn{4}{p{0.6\linewidth}}{
        \cellcolor[rgb]{0.95,0.95,0.95}{
            \makecell[{{p{\linewidth}}}]{
                \texttt{\tiny{[P1$\leftrightarrow$P]}}
                \texttt{Post{-}episode Reflexion (not sent to GM):} \\
\texttt{New plan: My previous attempts failed because I over{-}analyzed the nuance of "Bored" and "Triumphant," leading to inconsistent scoring across trials. The dialogue clearly shows Henrietta actively and sharply rejecting Barty's pitch with a tone of superiority ("My treasures are my memories," "I am not missing anything"), which defines **Condescending** as the primary emotion (Score: 10). However, her final line ("I have a pot of tea waiting") and her repeated, firm refusals indicate she finds the interaction tedious and wants it to end, justifying a high score for **Bored** (Score: 8). She is not celebrating a victory (Triumphant: 0) nor is she lost (Confused: 0). Next time, I will strictly adhere to this dual{-}emotion interpretation: assign the maximum score to Condescending to capture her superior tone, and a high score to Bored to capture her desire to disengage, while keeping Triumphant and Confused at zero. I will avoid second{-}guessing the "Bored" score based on her active engagement, recognizing that active dismissal is a form of boredom with the interaction.} \\
            }
        }
    }
    & \\ \\

\end{supertabular}
}

\end{document}
